\section{The Einstein Equation as a Fixed-Point Structure}
\label{sec:einstein}

\subsection{Positioning the Question}

So far: distance is derived from $C_{ij}$;
gravity appears as closure density $L$;
the metric takes the form $g_{\mu\nu} = g^{(0)}_{\mu\nu} + \alpha_G L_{\mu\nu}$.
The next question is therefore unavoidable:
\begin{equation}
  \boxed{\text{What equation does this structure obey?}}
\end{equation}
In conventional physics, the Einstein equation
$G_{\mu\nu} = 8\pi G T_{\mu\nu}$
is postulated as a fundamental law.
Within VED, the situation is reversed.
\begin{equation}
  \boxed{\text{The gravitational equation is not assumed.}}
\end{equation}
Instead, it must emerge from the same self-consistent closure structure.

\subsection{Starting Point: The Fixed-Point Equation}

Vol.~1 introduced the self-referential structure:
\begin{equation}
  C_{ij} = 1 - \exp(-\lambda \rho_i H_{ij}),
  \qquad
  \rho_i = \sum_j C_{ij}.
\end{equation}
This defines a nonlinear fixed-point problem.
\begin{equation}
  \boxed{\text{Structure determines itself.}}
\end{equation}
Geometry, density, and flow are all determined simultaneously
by this self-consistency condition.
This self-referential nature is the key to the emergence of gravitational
field equations.

\subsection{Continuum Limit and Vorticity}

Upon coarse-graining, define the flow $J_{ij} = C_{ij} - C_{ji}$.
In the continuum limit:
\begin{equation}
  \omega = \nabla \times J, \qquad L \sim |\omega|.
\end{equation}
\begin{equation}
  \boxed{\text{Closure density appears as vorticity of the flow.}}
\end{equation}
This connects the microscopic closure structure to macroscopic geometric distortion.

\subsection{Fixed-Point Condition}

In the classical phase established in Vol.~1:
\begin{equation}
  \sigma \approx 0, \qquad \nabla \cdot J \approx 0.
\end{equation}
These conditions correspond to approximate conservation.
\begin{equation}
  \boxed{\text{Conservation laws are properties of the phase.}}
\end{equation}
Under these conditions, closure structure stabilizes and the metric becomes
time-independent. This is the regime in which gravitational geometry emerges.

\subsection{Self-Consistency Condition for the Metric}

From Section~\ref{sec:closure_density},
$g_{\mu\nu} = g^{(0)}_{\mu\nu} + \alpha_G L_{\mu\nu}$.
However, $L_{\mu\nu}$ depends on flow, which depends on distance,
which depends on the metric itself.
Therefore $L_{\mu\nu} = L_{\mu\nu}[g]$, and the metric must satisfy:
\begin{equation}
  \label{eq:fixed_point_metric}
  \boxed{g_{\mu\nu} = g^{(0)}_{\mu\nu} + \alpha_G L_{\mu\nu}[g].}
\end{equation}
This defines a fixed-point problem for the metric.

\subsection{Linearization and the Form of the Equation}

In the weak-gravity limit $g_{\mu\nu} = \eta_{\mu\nu} + h_{\mu\nu}$
($|h_{\mu\nu}| \ll 1$),
expanding equation~\eqref{eq:fixed_point_metric} to first order yields:
\begin{equation}
  \label{eq:einstein_analog}
  \boxed{G_{\mu\nu} = 8\pi G_{\text{eff}} L_{\mu\nu}.}
\end{equation}
\begin{equation}
  \boxed{\text{Structurally analogous to the Einstein equation.}}
\end{equation}
Thus the Einstein equation appears as the linear approximation
of the self-consistent closure structure.

\subsection{The Inversion}

Conventional understanding: the Einstein equation = a fundamental law.

VED:
\begin{equation}
  \boxed{
    \text{The Einstein equation is a linear approximation of a fixed-point structure.}
  }
\end{equation}
It holds in the classical phase where closure is stable,
but need not hold outside this regime.
Gravity is therefore not fundamental but emergent.

\subsection{Why General Relativity Works}

This interpretation explains the empirical success of GR.
\begin{equation}
  \boxed{\text{GR is not incorrect.}}
\end{equation}
\begin{equation}
  \boxed{\text{It is an effective theory of the classical closure phase.}}
\end{equation}
In this phase, closure density varies slowly, conservation approximately holds,
and the metric is stable.
Under these conditions, the Einstein equation provides an accurate description.
